\conclusion

В ходе выполнения данной работы был разработан кроссплатформенный программный комплекс для тестирования и сравнения эффективности алгоритмов управления мульти-роботизированными системами. Поставленная цель достигнута в полном объёме, все сформулированные задачи решены.

В рамках работы выполнено следующее:

\begin{enumerate}[arabic]
    \item проведён анализ существующих алгоритмов решения задач MAPF и MRTA, выполнена их классификация по ключевым критериям: архитектуре управления, оптимальности и режиму работы;
    \item разработана модульная архитектура программного комплекса на основе паттерна <<Стратегия>>, обеспечивающая расширяемость --- добавление новых алгоритмов без модификации существующего кода;
    \item реализованы восемь MAPF-алгоритмов (CBS, CBS-SIPP, ECBS, EECBS, M*, CA*, WHCA*, PIBT), охватывающих оптимальные, субоптимальные и эвристические подходы;
    \item реализованы пять MRTA-алгоритмов (венгерский, жадный, последовательный аукцион, минимальная стоимость потока, комбинаторный аукцион) в двух режимах назначения: online и queue;
    \item разработан графический интерфейс, включающий редактор карт с возможностью сохранения и загрузки сценариев в формате JSON, экран выбора алгоритмов и визуализацию симуляции с отображением метрик в реальном времени;
    \item реализован режим пакетного тестирования с автоматической генерацией сценариев, четырьмя стратегиями размещения агентов, контролем таймаутов и экспортом результатов в CSV и графики;
    \item обеспечена кроссплатформенность: приложение собирается в автономный исполняемый файл с помощью PyInstaller для Windows, Linux и macOS;
    \item проведён сравнительный анализ реализованных алгоритмов, выявивший характерные области применимости каждого из них.
\end{enumerate}

Сравнительный анализ показал, что оптимальные алгоритмы (CBS, M*) обеспечивают наилучшее качество решений при малом числе агентов, субоптимальные (ECBS, EECBS) --- оптимальный баланс скорости и качества для средних сценариев, а децентрализованные (CA*, WHCA*, PIBT) --- наилучшую масштабируемость при большом числе агентов. Разработанный программный комплекс позволяет воспроизводить подобный анализ для произвольных конфигураций среды, что делает его полезным инструментом для обоснованного выбора алгоритмов в конкретных прикладных задачах.
